\begin{definition}[Subformula property]
A propositional derivation has the subformula property iff every
formula occurrence in it is a subformula of a formula in its endsequent.
\end{definition}

\begin{lemma}[Subformula property of cut-free derivations]
\label{lem:cutfree-subformula}
Every cut-free propositional \(\LK\)- or \(\LJ\)-derivation has the
subformula property.
\textit{Status: complete proof.}
\end{lemma}

\begin{proof}
In a cut-free derivation, inspect each rule upward from the endsequent.  Every
premise formula is either copied from the conclusion or is an immediate
subformula of the principal formula.  Induction on derivation height completes
the argument.
\end{proof}

\begin{corollary}[Propositional subformula property]
Every derivable propositional \(\LK\)- or \(\LJ\)-sequent has a cut-free
derivation with the subformula property.
\textit{Status: proof sketch.}
\end{corollary}

\begin{proofplan}
The cut-elimination theorem transforms a given derivation of the sequent
into a cut-free derivation of the same endsequent, and
\cref{lem:cutfree-subformula} shows that every such derivation has the
subformula property.
This derivation uses the preceding proof sketch of the cut-elimination
theorem; a complete proof supplies its reduction case table.
\end{proofplan}

\begin{remark}[First-order qualification]
For first-order cut-free derivations, quantifier rules introduce substitution
instances such as \(A[t/x]\).  The literal propositional subformula property
therefore fails.  The correct statement is a subformula property modulo term
substitution and eigenvariable renaming.  State the selected formulation
before using it in proof search.
\end{remark}

\begin{corollary}[Underivability of the empty sequent]
There is no \(\LK\)-derivation of the empty sequent
\(\varnothing\Rightarrow\varnothing\).
\textit{Status: proof sketch.}
\end{corollary}

\begin{proofplan}
If the empty sequent were derivable, cut elimination would give a cut-free
derivation of it.  The empty sequent is not an instance of a zero-premise
rule, and the conclusion of every other cut-free rule contains at least one
displayed formula occurrence.  Hence no cut-free derivation has the empty
sequent at its root, a contradiction.
This derivation uses the preceding proof sketch of the cut-elimination
theorem; a complete proof supplies its reduction case table.
\end{proofplan}

